\chapter{Integral Calculus in Multidimensional Spaces}

\section{Double Integrals and Planar Regions}
\label{sec:double-integrals}

The Riemann integral, as developed for functions of a single real variable, measures the signed area between the graph of a function and the interval on which it is defined, via a limiting process applied to sums that approximate this area by finitely many rectangles. The same idea extends, with only notational complications, to functions of several variables: a function of two variables may be integrated over a planar region to yield a signed volume, a function of three variables integrated over a solid region yields a signed hypervolume, and the construction generalizes without further conceptual difficulty to functions on $\mathbb{R}^n$. This chapter develops the theory of multiple integration from its definition via partitions and Darboux sums, through Fubini's Theorem, which reduces the computation of a multiple integral to a sequence of ordinary single-variable integrations, to the change of variables formula, which generalizes $u$-substitution to several dimensions and relies essentially on the Jacobian matrix introduced in the preceding chapter.

\subsection{Construction of the Double Integral on Rectangles}

As in the one-dimensional theory, the integral is defined first over the simplest possible domain, a rectangle, and is subsequently extended to more general regions.

\begin{definition}[Partition of a Rectangle]
Consider the rectangle
\[
R = [a,b]\times[c,d] \subset \mathbb{R}^2.
\]
A partition of $R$ is a pair $P = (P_x, P_y)$, where
\[
\begin{gathered}
P_x : a = x_0 < x_1 < \cdots < x_m = b \\
P_y : c = y_0 < y_1 < \cdots < y_n = d
\end{gathered}
\]
are partitions of $[a,b]$ and $[c,d]$ respectively. The partition $P$ divides $R$ into $mn$ closed subrectangles
\[
R_{ij} = [x_{i-1},x_i]\times[y_{j-1},y_j],
\]
of area
\[
\Delta A_{ij} = \Delta x_i \, \Delta y_j = (x_i - x_{i-1})(y_j-y_{j-1}).
\]
The mesh of $P$ is
\[
\|P\| = \max_{i,j} \max(\Delta x_i, \Delta y_j).
\]

\end{definition}

\begin{definition}[Darboux Sums and the Double Integral]
Let $f : R \to \mathbb{R}$ be bounded, and let $P$ be a partition of $R$ with subrectangles $R_{ij}$. Set
\[
M_{ij} = \sup_{(x,y)\in R_{ij}} f(x,y), \qquad m_{ij} = \inf_{(x,y)\in R_{ij}} f(x,y),
\]
and define the upper and lower Darboux sums
\[
U(f,P) = \sum_{i,j} M_{ij}\,\Delta A_{ij}, \qquad L(f,P) = \sum_{i,j} m_{ij}\,\Delta A_{ij}.
\]
The function $f$ is (Darboux) integrable on $R$ if
\[
\sup_P L(f,P) = \inf_P U(f,P),
\]

In this case, the common value is called the double integral of $f$ over $R$. The following two notations are equivalent:
\[
\begin{gathered}
\iint_R f \, dA, \\
\iint_R f(x,y)\, dx\, dy.
\end{gathered}
\]

\end{definition}

Since refining a partition can only decrease upper sums and increase lower sums, one has
\[
L(f,P) \leq \sup_{P'} L(f,P') \leq \inf_{P'} U(f,P') \leq U(f,P)
\]
for every partition $P$, exactly as in the one-dimensional theory; the content of integrability is that these two bounding quantities coincide.

\begin{theorem}[Continuity Implies Integrability]
If $f$ is continuous on the closed rectangle
\[
R = [a,b]\times[c,d],
\]
then $f$ is integrable on $R$.
\end{theorem}

\begin{proof}
Since $R$ is closed and bounded, it is compact, and since $f$ is continuous on $R$, it is uniformly continuous on $R$ by the Heine--Cantor theorem: given $\varepsilon>0$, there exists $\delta>0$ such that $\|(x,y)-(x',y')\| < \delta$ implies
\[
|f(x,y)-f(x',y')| < \varepsilon / \operatorname{Area}(R).
\]
Choose any partition $P$ with mesh
\[
\|P\| < \delta/\sqrt{2},
\]
so that every subrectangle $R_{ij}$ has diameter
\[
((\Delta x_i)^2 + (\Delta y_j)^2)^{1/2} < \delta.
\]

On each $R_{ij}$, which is compact, $f$ attains its supremum $M_{ij}$ and infimum $m_{ij}$ at some points $\mathbf{p}_{ij}, \mathbf{q}_{ij} \in R_{ij}$ by the Extreme Value Theorem; since $\|\mathbf{p}_{ij}-\mathbf{q}_{ij}\| < \delta$, uniform continuity gives
\[
M_{ij}-m_{ij} = f(\mathbf{p}_{ij})-f(\mathbf{q}_{ij}) < \varepsilon/\operatorname{Area}(R).
\]
Consequently
\[
U(f,P) - L(f,P) = \sum_{i,j}(M_{ij}-m_{ij})\,\Delta A_{ij} < \frac{\varepsilon}{\operatorname{Area}(R)}\sum_{i,j}\Delta A_{ij} = \varepsilon,
\]
since
\[
\sum_{i,j}\Delta A_{ij} = \operatorname{Area}(R).
\]
For every partition, the following two inequalities hold:
\[
\begin{gathered}
\sup_{P'}L(f,P') \geq L(f,P) \\
\inf_{P'}U(f,P') \leq U(f,P)
\end{gathered}
\]
For the particular partition $P$ chosen above, they imply
\[
0 \leq \inf_{P'}U(f,P') - \sup_{P'}L(f,P') \leq U(f,P)-L(f,P) < \varepsilon
\]
Since $\varepsilon>0$ was arbitrary, it follows that
\[
\inf_{P'}U(f,P') = \sup_{P'}L(f,P'),
\]
that is, $f$ is integrable on $R$.
\end{proof}

Fubini's Theorem is the principal computational tool of multivariable integration: it asserts that, under mild hypotheses, a double integral may be evaluated as an iterated sequence of two ordinary one-dimensional integrals, performed one variable at a time.

\begin{theorem}[Fubini's Theorem, Rectangular Case]
Let $f$ be continuous on
\[
R = [a,b]\times[c,d].
\]

Then, for every fixed $x \in [a,b]$, the function $y \mapsto f(x,y)$ is continuous on $[c,d]$ and hence integrable there, and the resulting function
\[
A(x) = \int_c^d f(x,y)\, dy
\]
is continuous, hence integrable, on $[a,b]$. Moreover,
\[
\iint_R f\, dA = \int_a^b A(x)\, dx = \int_a^b \left( \int_c^d f(x,y)\, dy \right) dx.
\]
The same conclusion holds with the roles of $x$ and $y$ exchanged, so that in particular
\[
\int_a^b \left( \int_c^d f(x,y)\, dy \right) dx = \int_c^d \left( \int_a^b f(x,y)\, dx \right) dy.
\]

\end{theorem}

\begin{proof}
Fix $x \in [a,b]$; since $f$ is continuous on $R$, the single-variable function $y \mapsto f(x,y)$ is continuous on $[c,d]$, and therefore integrable there by the one-dimensional theory, so $A(x)$ is well defined. Let $P = (P_x,P_y)$ be any partition of $R$, with subrectangles
\[
R_{ij} = [x_{i-1},x_i]\times[y_{j-1},y_j],
\]
and retain the notation $M_{ij}, m_{ij}$ of the preceding section. Fix $i$, and let $x \in [x_{i-1},x_i]$. For every $j$ and every $y \in [y_{j-1},y_j]$, one has
\[
m_{ij} \leq f(x,y) \leq M_{ij},
\]
since $(x,y) \in R_{ij}$; integrating this inequality over $y \in [y_{j-1},y_j]$ gives
\[
m_{ij}\,\Delta y_j \;\leq\; \int_{y_{j-1}}^{y_j} f(x,y)\, dy \;\leq\; M_{ij}\,\Delta y_j.
\]

Summing over $j = 1, \dots, n$, and using additivity of the one-dimensional integral over the subintervals of $[c,d]$,
\[
\sum_{j} m_{ij}\,\Delta y_j \;\leq\; A(x) = \int_c^d f(x,y)\, dy \;\leq\; \sum_j M_{ij}\, \Delta y_j,
\]
valid for every $x \in [x_{i-1},x_i]$. Integrating this inequality over $x \in [x_{i-1},x_i]$ and summing over $i = 1, \dots, m$ gives
\[
\sum_{i,j} m_{ij}\, \Delta x_i \, \Delta y_j \;\leq\; \int_a^b A(x)\, dx \;\leq\; \sum_{i,j} M_{ij}\,\Delta x_i\, \Delta y_j,
\]
For every partition $P$ of $R$, this gives the bounds
\[
L(f,P) \;\leq\; \int_a^b A(x)\, dx \;\leq\; U(f,P),
\]
Since $f$ is continuous on $R$, it is integrable there by the previous theorem, so
\[
\sup_P L(f,P) = \inf_P U(f,P) = \iint_R f\, dA;
\]
as the fixed number
\[
\int_a^b A(x)\, dx
\]
is squeezed between $L(f,P)$ and $U(f,P)$ for every partition $P$, it must equal this common value, giving
\[
\int_a^b A(x)\, dx = \iint_R f\, dA.
\]

Continuity of $A$ follows from uniform continuity of $f$ on the compact rectangle $R$: given $\varepsilon>0$, choose $\delta>0$ so that $\|(x,y)-(x',y)\|<\delta$ implies
\[
|f(x,y)-f(x',y)| < \varepsilon/(d-c)
\]
for all $y$; then $|x-x'|<\delta$ implies
\[
|A(x)-A(x')| \leq \int_c^d |f(x,y)-f(x',y)|\, dy < \varepsilon,
\]
establishing continuity, hence integrability, of $A$. The identity with the roles of $x$ and $y$ exchanged follows by the identical argument applied with the two coordinates interchanged throughout.
\end{proof}

\begin{example}
Let $R = [0,1]\times[0,2]$ and $f(x,y) = xy^2$. By Fubini's Theorem,
\begin{align*}
\iint_R xy^2\,\dd A
&= \int_0^1 \left(\int_0^2 xy^2\,\dd y\right)\dd x \\
&= \int_0^1 x\left[\frac{y^3}{3}\right]_{0}^{2}\dd x \\
&= \int_0^1 \frac{8x}{3}\,\dd x
= \frac{8}{3}\left[\frac{x^2}{2}\right]_{0}^{1}
= \frac{4}{3}.
\end{align*}
Integrating in the opposite order gives the same value:
\begin{align*}
\int_0^2\left(\int_0^1 xy^2\,\dd x\right)\dd y
&= \int_0^2 \frac{y^2}{2}\,\dd y \\
&= \left[\frac{y^3}{6}\right]_{0}^{2}
= \frac{4}{3},
\end{align*}
confirming the theorem.
\end{example}

\subsection{Integrability on General Planar Regions}


Most regions of interest are not rectangles, and the integral must be extended to accommodate them.

\begin{definition}[Content Zero]
A bounded set $Z \subset \mathbb{R}^2$ has content zero if, for every $\varepsilon>0$, there exist finitely many rectangles $S_1,\dots,S_k$ with
\[
\begin{gathered}
Z \subseteq \bigcup_{l=1}^{k} S_l \\
\sum_{l=1}^{k}\operatorname{Area}(S_l) < \varepsilon.
\end{gathered}
\]

\end{definition}

The graph of a continuous function on a closed bounded interval has content zero, a fact used implicitly below without separate proof, since it follows from uniform continuity in the same manner as the theorem above: covering $[a,b]$ by finitely many subintervals of total length $\varepsilon'$ and small enough that the corresponding oscillation of the function is at most $\varepsilon/(b-a)$, the rectangles built from these subintervals and this oscillation bound cover the graph with total area at most $\varepsilon$.

\begin{theorem}[Integrability with a Content-Zero Set of Discontinuities]
Let $\tilde{f} : R \to \mathbb{R}$ be bounded on a rectangle $R$, say $|\tilde{f}| \leq M$, and continuous at every point of $R$ except possibly on a set $Z \subset R$ of content zero. Then $\tilde{f}$ is integrable on $R$.
\end{theorem}

\begin{proof}
Let $\varepsilon > 0$. Since $Z$ has content zero, cover $Z$ by finitely many open rectangles $S_1,\dots,S_k$ with
\[
\sum_l \operatorname{Area}(S_l) < \varepsilon/(4M).
\]

The set $K = R \setminus \bigcup_l S_l$ is closed and bounded, hence compact, and $\tilde{f}$ is continuous at every point of $K$, since $K$ contains no point of $Z$; $\tilde{f}$ is therefore uniformly continuous on $K$. Choose a partition $P$ of $R$ fine enough that: (i) every subrectangle of $P$ meeting $\bigcup_l S_l$ is contained in a slightly enlarged union of the $S_l$ whose total area remains below $\varepsilon/(2M)$, which is possible since finitely many rectangles have a boundary of content zero and the $S_l$ may be taken slightly larger at the outset to absorb this adjustment; and (ii) the mesh of $P$ is small enough that, on every subrectangle disjoint from $\bigcup_l S_l$ (hence contained in $K$), the oscillation of $\tilde f$ is less than $\varepsilon/(2\operatorname{Area}(R))$, which is possible by uniform continuity of $\tilde f$ on $K$. Splitting the sum defining $U(f,P)-L(f,P)$ into subrectangles of the first kind and of the second kind,
\[
U(\tilde f,P) - L(\tilde f,P) = \sum_{\text{meets } \cup S_l} (M_{ij}-m_{ij})\Delta A_{ij} \;+\; \sum_{\text{in } K} (M_{ij}-m_{ij}) \Delta A_{ij}.
\]

In the first sum, $M_{ij}-m_{ij} \leq 2M$ always, so this sum is at most $2M \cdot \varepsilon/(4M) = \varepsilon/2$; in the second sum, $M_{ij}-m_{ij} < \varepsilon/(2\operatorname{Area}(R))$, so this sum is less than
\[
\varepsilon/(2\operatorname{Area}(R)) \cdot \operatorname{Area}(R) = \varepsilon/2.
\]

Hence $U(\tilde f,P)-L(\tilde f,P) < \varepsilon$; since $\varepsilon>0$ was arbitrary, $\tilde f$ is integrable on $R$, exactly as in the proof of the theorem for continuous functions above.
\end{proof}

\begin{definition}[Integral over a General Bounded Region]
Let $D \subset \mathbb{R}^2$ be bounded, with $\partial D$ of content zero, and let $f$ be bounded on $D$ and continuous except possibly on a set of content zero. Choose any rectangle $R \supseteq D$, and extend $f$ to $\tilde f$ on $R$ by setting $\tilde f = f$ on $D$ and $\tilde f = 0$ on $R \setminus D$. The double integral of $f$ over $D$ is defined as
\[
\iint_D f\, dA := \iint_R \tilde f\, dA.
\]

\end{definition}

The discontinuities of $\tilde f$ occur only within $\partial D$, together with the discontinuities of $f$ already present in $D$, together forming a set of content zero; the preceding theorem guarantees that $\tilde f$ is integrable on $R$, so this definition is meaningful, and it is independent of the choice of enclosing rectangle $R$, since $\tilde f$ vanishes outside $D$ regardless.

\begin{definition}[Type I and Type II Regions]
A region $D \subset \mathbb{R}^2$ is of Type I if
\[
D = \{(x,y) : a \leq x \leq b, \; g_1(x) \leq y \leq g_2(x)\}
\]
for continuous functions $g_1 \leq g_2$ on $[a,b]$; it is of Type II if
\[
D = \{(x,y) : c \leq y \leq d, \; h_1(y) \leq x \leq h_2(y)\}
\]
for continuous functions $h_1 \leq h_2$ on $[c,d]$.
\end{definition}

For a Type I region, $\partial D$ consists of the graphs of $g_1$ and $g_2$, each of content zero as noted above, together with the two vertical segments $x=a$ and $x=b$, also visibly of content zero; hence $\partial D$ has content zero, being a finite union of sets of content zero, and the integral over $D$ is defined. Applying Fubini's Theorem on an enclosing rectangle $R = [a,b]\times[c,d]$ to the extension $\tilde f$, and noting that $\tilde f(x,y) = 0$ for $y \notin [g_1(x),g_2(x)]$, gives the following practical evaluation formula.

\begin{theorem}[Iterated Integral over a Type I Region]
If $D$ is a Type I region as above and $f$ is continuous on $D$, then
\[
\iint_D f\, dA = \int_a^b \left( \int_{g_1(x)}^{g_2(x)} f(x,y)\, dy \right) dx.
\]
An analogous formula, with the roles of $x$ and $y$ exchanged, holds for a Type II region.
\end{theorem}

\begin{proof}
Let $R = [a,b]\times[c,d]$ enclose $D$, with $c \leq g_1(x)$ and $g_2(x) \leq d$ for all $x \in [a,b]$. By Fubini's Theorem applied to $\tilde f$ on $R$,
\[
\iint_D f \, dA = \iint_R \tilde f \, dA = \int_a^b \left( \int_c^d \tilde f(x,y)\, dy \right) dx.
\]

For fixed $x \in [a,b]$, $\tilde f(x,y) = f(x,y)$ for $y \in [g_1(x),g_2(x)]$ and $\tilde f(x,y) = 0$ otherwise, so by additivity of the one-dimensional integral over subintervals,
\[
\int_c^d \tilde f(x,y)\, dy = \int_c^{g_1(x)} 0 \, dy + \int_{g_1(x)}^{g_2(x)} f(x,y)\, dy + \int_{g_2(x)}^d 0\, dy = \int_{g_1(x)}^{g_2(x)} f(x,y)\, dy,
\]
which, substituted into the displayed identity above, gives the claimed formula.
\end{proof}

\begin{theorem}[Basic Properties of the Double Integral]
Let $f,g$ be integrable on a region $D$, and let $\alpha,\beta \in \mathbb{R}$. Then:
\begin{enumerate}
\item[(i)] \textit{Linearity}: the following identity holds:

\[
\iint_D (\alpha f + \beta g)\, dA = \alpha \iint_D f\, dA + \beta \iint_D g\, dA.
\]

\item[(ii)] \textit{Monotonicity}: if $f \leq g$ on $D$, then

\[
\iint_D f\, dA \leq \iint_D g\, dA.
\]

\item[(iii)] \textit{Additivity over regions}: if $D = D_1 \cup D_2$ with $D_1, D_2$ overlapping at most on a set of content zero, then

\[
\iint_D f\, dA = \iint_{D_1} f\, dA + \iint_{D_2} f\, dA.
\]

\end{enumerate}
\end{theorem}

\begin{proof}
(i) and (ii) follow directly from the corresponding properties of upper and lower Darboux sums: since
\[
M_{ij}[\alpha f + \beta g] = \alpha M_{ij}[f] + \beta M_{ij}[g]
\]
whenever $\alpha,\beta \geq 0$ (with the analogous relation for $m_{ij}$ when either scalar is negative, upon exchanging the roles of $M_{ij}$ and $m_{ij}$), the Darboux sums of $\alpha f + \beta g$ are the corresponding linear combinations of the Darboux sums of $f$ and $g$ for every partition, and passing to the common limiting value defining the integral yields (i); (ii) follows since $f \leq g$ on $D$ forces $L(f,P) \leq L(g,P)$ for every partition $P$ of an enclosing rectangle (applied to the zero-extensions), hence
\[
\iint_D f\,dA = \sup_P L(f,P) \leq \sup_P L(g,P) = \iint_D g\, dA.
\]

For (iii), extend $f$ by zero to a common enclosing rectangle $R$; the zero-extension of $f$ from $D$ agrees with the sum of the zero-extensions from $D_1$ and from $D_2$ except on $D_1 \cap D_2$, a set of content zero which therefore does not affect any Darboux sum in the limit, by the same estimate used in the proof of the integrability theorem above; hence
\[
\iint_D f \, dA = \iint_{D_1} f\, dA + \iint_{D_2} f \, dA
\]
by linearity applied to the (almost everywhere) sum of the two zero-extensions.
\end{proof}

Figure~\ref{fig:ch11-typeI-region} shows the vertical slices used to evaluate an integral over a Type I region:
\[
\iint_D f\,dA.
\]

\begin{figure}[!htbp]
\centering
\begin{tikzpicture}[scale=1.15,>=Stealth]
  \draw[->] (-0.4,0) -- (4.4,0) node[right,font=\scriptsize]{$x$};
  \draw[->] (0,-0.4) -- (0,3.3) node[above,font=\scriptsize]{$y$};
  \fill[figblue,opacity=0.30]
    plot[domain=0.6:3.6,smooth,variable=\x] ({\x},{0.35+0.18*(\x-0.6)*(\x-0.6)})
    -- plot[domain=3.6:0.6,smooth,variable=\x] ({\x},{2.9-0.28*(\x-2.1)*(\x-2.1)})
    -- cycle;
  \draw[vecB,thick,domain=0.6:3.6,smooth,variable=\x]
    plot ({\x},{0.35+0.18*(\x-0.6)*(\x-0.6)}) node[right,font=\scriptsize]{$y=g_1(x)$};
  \draw[vecC,thick,domain=0.6:3.6,smooth,variable=\x]
    plot ({\x},{2.9-0.28*(\x-2.1)*(\x-2.1)}) node[right,font=\scriptsize]{$y=g_2(x)$};
  \draw[dashed,gray] (0.6,0) -- (0.6,0.35) node[pos=0,below,font=\scriptsize,black]{$a$};
  \draw[dashed,gray] (3.6,0) -- (3.6,1.62);
  \node[below,font=\scriptsize] at (3.6,0) {$b$};
  \node[font=\scriptsize] at (2.0,1.5) {$D$};
\end{tikzpicture}
\caption[Type I region]{A Type I region bounded by two graphs and two vertical lines. Integrating along each vertical slice first, and then across the horizontal interval, gives the double integral over $D$.}
\label{fig:ch11-typeI-region}
\end{figure}

Figure~\ref{fig:atlas-triangular-domain} compares vertical and horizontal slices of the same integration region.

\begin{figure}[!htbp]
\centering
\begin{tikzpicture}[>=Stealth]
\begin{scope}[scale=3.6]
\fill[figblue!15] (0,0) -- (1,0) -- (0,1) -- cycle;
\draw[->] (-0.1,0) -- (1.2,0) node[right] {$x$};
\draw[->] (0,-0.1) -- (0,1.2) node[above] {$y$};
\draw[figblue,thick] (0,1) -- (1,0);
\draw[figred,very thick] (0.35,0) -- (0.35,0.65);
\node[below] at (0.35,0) {$x$};
\node[above right] at (0.35,0.65) {$1-x$};
\node[below] at (1,0) {$1$};
\node[left] at (0,1) {$1$};
\end{scope}
\begin{scope}[shift={(6,0)},scale=3.6]
\fill[figblue!15] (0,0) -- (1,0) -- (0,1) -- cycle;
\draw[->] (-0.1,0) -- (1.2,0) node[right] {$x$};
\draw[->] (0,-0.1) -- (0,1.2) node[above] {$y$};
\draw[figblue,thick] (0,1) -- (1,0);
\draw[figgreen,very thick] (0,0.4) -- (0.6,0.4);
\node[left] at (0,0.4) {$y$};
\node[above right] at (0.6,0.4) {$1-y$};
\node[below] at (1,0) {$1$};
\node[left] at (0,1) {$1$};
\end{scope}
\end{tikzpicture}
\caption{The triangular region is bounded by the coordinate axes and $x+y=1$. Vertical slices use $0\le y\le1-x$, whereas horizontal slices use $0\le x\le1-y$. Both descriptions cover the same region and lead to equivalent iterated integrals.}
\label{fig:atlas-triangular-domain}
\end{figure}

\subsection{Change of Variables in the Plane}

The one-dimensional substitution rule
\[
\int_a^b f(x)\,dx = \int_{\alpha}^{\beta} f(\varphi(u))\,\varphi'(u)\,du,
\]
valid for a differentiable bijection $\varphi$ with $\varphi(\alpha)=a$, $\varphi(\beta)=b$, admits a direct generalization to $\mathbb{R}^2$, in which the scalar factor $\varphi'(u)$ is replaced by the absolute value of the Jacobian determinant introduced in the preceding chapter.

\begin{theorem}[Change of Variables Formula, Two Dimensions]
Let $T : U \to V$,
\[
T(u,v) = (x(u,v),y(u,v)),
\]
be a bijective $C^1$ map between open sets $U,V \subset \mathbb{R}^2$, with $\det J_T(u,v) \neq 0$ for every $(u,v) \in U$. Let $D \subset U$ be a bounded region with $D' = T(D) \subset V$, and let $f$ be continuous on $D'$. Then
\[
\iint_{D'} f(x,y)\, dx\, dy = \iint_D f\big(T(u,v)\big)\, |\det J_T(u,v)|\, du\, dv.
\]

\end{theorem}

\begin{proof}[Derivation]
Partition $D$ by a fine grid of small rectangles
\[
R_{ij} = [u_{i-1},u_i]\times[v_{j-1},v_j],
\]
of dimensions $\Delta u \times \Delta v$, and fix a corner $(u_i,v_j)$ of $R_{ij}$. By differentiability of $T$ at $(u_i,v_j)$ -- precisely the statement that $J_T(u_i,v_j)$ furnishes the best local linear approximation to $T$, established in the preceding chapter -- the image of $R_{ij}$ under $T$ satisfies
\begin{align*}
T(u_i + \Delta u, v_j) &\approx T(u_i,v_j) + \frac{\partial T}{\partial u}(u_i,v_j)\,\Delta u, \\
T(u_i, v_j+\Delta v) &\approx T(u_i,v_j) + \frac{\partial T}{\partial v}(u_i,v_j)\, \Delta v,
\end{align*}
with error terms that are $o(\Delta u)$ and $o(\Delta v)$ respectively as $\Delta u, \Delta v \to 0$. Consequently, $T(R_{ij})$ is approximated, to leading order, by the parallelogram based at $T(u_i,v_j)$ and spanned by the two vectors $\partial T/\partial u\, \Delta u$ and $\partial T/\partial v\, \Delta v$, which are precisely the two columns of $J_T(u_i,v_j)$, each scaled respectively by $\Delta u$ and $\Delta v$. The area of a parallelogram spanned by two vectors equals the absolute value of the determinant of the matrix formed by those vectors as columns; hence
\begin{align*}
\operatorname{Area}\big(T(R_{ij})\big) &\;\approx\; \left| \det \begin{pmatrix} \partial T/\partial u \, \Delta u & \partial T/\partial v \, \Delta v \end{pmatrix} \right| \\
&= |\det J_T(u_i,v_j)| \, \Delta u\, \Delta v = |\det J_T(u_i,v_j)|\, \operatorname{Area}(R_{ij}),
\end{align*}
using multilinearity of the determinant to factor $\Delta u$ and $\Delta v$ out of the respective columns. As the mesh of the partition tends to zero, the images $T(R_{ij})$ tile $D'$ with overlaps and gaps of area tending to zero relative to $\operatorname{Area}(D')$, by injectivity and continuity of $T$, so that
\begin{align*}
\iint_{D'} f\, dA &= \lim_{\text{mesh}\to 0} \sum_{i,j} f\big(T(u_i,v_j)\big)\, \operatorname{Area}\big(T(R_{ij})\big) \\
&\;\approx\; \lim_{\text{mesh}\to 0} \sum_{i,j} f\big(T(u_i,v_j)\big)\, |\det J_T(u_i,v_j)|\, \Delta u \, \Delta v,
\end{align*}
and the right-hand side is precisely a Riemann sum, over the partition $\{R_{ij}\}$ of $D$, for the integrand $f(T(u,v))\,|\det J_T(u,v)|$, which is continuous on $D$ since $f$, $T$, and $\det J_T$ are all continuous, hence integrable there. Passing to the limit converts both approximate equalities into exact ones, since the accumulated error, controlled by the uniform continuity of $J_T$ on the compact closure of $D$, vanishes in the limit; we omit the fully detailed uniform error estimates, which require only a routine, if tedious, application of uniform continuity of $DT$ on compact subsets. This establishes the formula.
\end{proof}

\begin{corollary}[Polar Coordinates]
Define the polar-coordinate map by
\[
T(r,\theta) = (r\cos\theta, r\sin\theta).
\]

As computed in the example of the preceding chapter, $\det J_T(r,\theta) = r$ (see Figure~\ref{fig:ch10-jacobian-polar}), so that for a region $D'$ described in polar coordinates by
\[
D = \{(r,\theta) : \alpha \leq \theta \leq \beta, \; r_1(\theta) \leq r \leq r_2(\theta)\},
\]

\[
\iint_{D'} f(x,y)\, dx\, dy = \int_{\alpha}^{\beta} \int_{r_1(\theta)}^{r_2(\theta)} f(r\cos\theta, r\sin\theta)\; r\, dr\, d\theta.
\]

\end{corollary}

\begin{example}[The Gaussian Integral]
The improper integral
\[
I = \int_{-\infty}^{\infty} e^{-x^2}\, dx
\]
admits no elementary antiderivative, yet its value can be found exactly by passing to two dimensions and then to polar coordinates. Since
\[
I = \int_{-\infty}^\infty e^{-y^2}\,dy
\]
as well, Fubini's Theorem (extended to the improper setting, valid here since the integrand is nonnegative) gives
\[
I^2 = \left(\int_{-\infty}^{\infty} e^{-x^2}\, dx\right)\left(\int_{-\infty}^{\infty} e^{-y^2}\, dy\right) = \iint_{\mathbb{R}^2} e^{-x^2-y^2}\, dA.
\]

Passing to polar coordinates, with $x^2+y^2 = r^2$ and the region $\mathbb{R}^2$ described by $r \in [0,\infty)$, $\theta \in [0,2\pi)$,
\[
\begin{aligned}
I^2 &= \int_0^{2\pi}\int_0^{\infty} e^{-r^2}\, r\, dr\, d\theta \\
&= 2\pi \int_0^{\infty} e^{-r^2} r\, dr \\
&= 2\pi \left[ -\frac{1}{2} e^{-r^2} \right]_0^{\infty}
= 2\pi \cdot \frac{1}{2} = \pi,
\end{aligned}
\]
where the substitution $t=r^2$, $dt = 2r\,dr$ converts
\[
\int_0^\infty e^{-r^2}r\,dr
\]
into the elementary
\[
\frac12\int_0^\infty e^{-t}\,dt = \frac12.
\]
Since $I > 0$, this gives the classical identity
\[
\int_{-\infty}^{\infty} e^{-x^2}\, dx = \sqrt{\pi}.
\]

\end{example}

\section{Triple Integrals and Spatial Coordinates}

\subsection{Triple Integrals over Boxes}

The construction of Section~\ref{sec:double-integrals} extends verbatim to three dimensions, replacing rectangles by boxes and areas by volumes.

\begin{definition}[Triple Integral over a Box]
Consider the box
\[
B = [a_1,b_1]\times[a_2,b_2]\times[a_3,b_3] \subset \mathbb{R}^3,
\]

Let $f : B \to \mathbb{R}$ be bounded. A partition of $B$ is obtained by partitioning each edge, dividing $B$ into subboxes $B_{ijk}$ of volume $\Delta V_{ijk}$. With
\[
\begin{gathered}
M_{ijk} = \sup_{B_{ijk}} f \\
m_{ijk} = \inf_{B_{ijk}} f,
\end{gathered}
\]
define
\[
\begin{gathered}
U(f,P) = \sum M_{ijk}\Delta V_{ijk} \\
L(f,P) = \sum m_{ijk} \Delta V_{ijk}
\end{gathered}
\]
exactly as before; $f$ is integrable on $B$ if
\[
\sup_P L(f,P) = \inf_P U(f,P),
\]
and this common value is the triple integral
\[
\iiint_B f\, dV.
\]

\end{definition}

The proofs that continuous functions on a box are integrable, and that integrability persists under a content-zero (now, volume-zero) set of discontinuities, are entirely analogous to the two-dimensional case, replacing the Heine--Cantor and covering arguments given above by their immediate three-variable counterparts; we do not repeat the details. The corresponding version of Fubini's Theorem likewise follows the same squeeze argument as before, applied with one variable singled out at a time.

\begin{theorem}[Fubini's Theorem, Box Case]
If $f$ is continuous on
\[
B = [a_1,b_1]\times[a_2,b_2]\times[a_3,b_3],
\]
Fubini's formula gives
\[
\iiint_B f\, dV = \int_{a_1}^{b_1}\int_{a_2}^{b_2}\int_{a_3}^{b_3} f(x,y,z)\; dz\, dy\, dx,
\]

The iterated integral may equally be evaluated in any of the six orders obtained by permuting $x,y,z$.
\end{theorem}

\begin{proof}
Group the last two variables together: writing
\[
\begin{gathered}
\mathbf{y} = (y,z) \in \mathbb{R}^2 \\
R = [a_2,b_2]\times[a_3,b_3],
\end{gathered}
\]
the two-dimensional Fubini Theorem, applied to the function $\mathbf{y} \mapsto f(x,\mathbf{y})$ for fixed $x$, together with the identical squeeze argument of the rectangular case applied with $x$ playing the role of the single ``outer'' variable and $R$ playing the role of the two-dimensional ``inner'' rectangle, yields
\[
\iiint_B f\, dV = \int_{a_1}^{b_1}\left(\iint_R f(x,y,z)\, dy\,dz\right) dx;
\]
applying the two-dimensional Fubini Theorem once more to the inner double integral, for each fixed $x$, resolves it into the iterated integral over $y$ then $z$ (or $z$ then $y$), giving the stated formula. Since the grouping of variables was arbitrary, the same argument applied with a different variable singled out as ``outer'' establishes the remaining orders of integration.
\end{proof}

\subsection{Triple Integrals over General Regions}

\begin{definition}[Type I Solid Region]
A solid region $E \subset \mathbb{R}^3$ is of Type I if
\[
E = \big\{ (x,y,z) : (x,y) \in D, \; u_1(x,y) \leq z \leq u_2(x,y) \big\}
\]
for a Type I or Type II planar region $D \subset \mathbb{R}^2$ and continuous functions $u_1 \leq u_2$ on $D$. Regions defined analogously with the roles of $x,y,z$ permuted are of Type II and Type III.
\end{definition}

As in the planar case, extending $f$ by zero outside $E$ and applying Fubini's Theorem on an enclosing box yields, for a Type I region,
\[
\iiint_E f\, dV = \iint_D \left( \int_{u_1(x,y)}^{u_2(x,y)} f(x,y,z)\, dz \right) dA,
\]
and the inner double integral over $D$ is then evaluated by the planar methods of the preceding sections.

\begin{example}[Volume of a Tetrahedron]
Let $E$ be the solid tetrahedron bounded by the coordinate planes $x=0$, $y=0$, $z=0$ and the plane $x+y+z=1$. Regarding $E$ as a Type I region over the planar triangle
\[
D = \{(x,y): x\geq0,\,y\geq0,\,x+y\leq1\},
\]
The vertical bounds over this base are
\[
\begin{aligned}
0 &\leq z \\
&\leq 1-x-y,
\end{aligned}
\]

\[
\operatorname{Vol}(E) = \iiint_E dV = \iint_D (1-x-y)\, dA = \int_0^1 \int_0^{1-x} (1-x-y)\, dy\, dx.
\]
The inner integral is
\[
\int_0^{1-x}(1-x-y)\,dy = \left[(1-x)y - \tfrac{y^2}{2}\right]_0^{1-x} = (1-x)^2 - \tfrac{(1-x)^2}{2} = \tfrac{(1-x)^2}{2},
\]
so
\[
\operatorname{Vol}(E) = \int_0^1 \frac{(1-x)^2}{2}\, dx = \left[-\frac{(1-x)^3}{6}\right]_0^1 = \frac{1}{6},
\]
matching the familiar formula
\[
\operatorname{Vol} = (1/6)|abc|
\]
for a coordinate tetrahedron with
\[
\begin{aligned}
a &= b \\
&= c \\
&= 1.
\end{aligned}
\]

\end{example}

\subsection{Cylindrical Coordinates}


The three-dimensional change of variables formula is the direct generalization of the planar one, with the parallelogram of the earlier derivation replaced by a parallelepiped, whose volume is likewise the absolute value of the determinant of the matrix formed by its three spanning vectors -- here, the three columns of the Jacobian matrix, scaled by the corresponding side lengths of the small box.

\begin{theorem}[Change of Variables Formula, Three Dimensions]
Let $T : U \to V$ be a bijective $C^1$ map between open sets $U,V \subset \mathbb{R}^3$ with $\det J_T \neq 0$ on $U$. Let $E \subset U$ with $E' = T(E) \subset V$, and let $f$ be continuous on $E'$. Then
\[
\iiint_{E'} f\, dV = \iiint_E f\big(T(u,v,w)\big)\, |\det J_T(u,v,w)| \; du\, dv\, dw.
\]

\end{theorem}

\begin{proof}[Derivation]
The argument is identical, mutatis mutandis, to the two-dimensional derivation above: partitioning $E$ into small boxes, approximating the image of each box under $T$ by the parallelepiped spanned by the three columns of $J_T$ scaled by $\Delta u, \Delta v, \Delta w$, whose volume is $|\det J_T|\,\Delta u\,\Delta v\,\Delta w$, and passing to the limit of Riemann sums; we do not repeat the details.
\end{proof}

Cylindrical coordinates $(r,\theta,z)$ describe a point via its polar coordinates $(r,\theta)$ in the $xy$-plane together with its height $z$, via the map
\[
T(r,\theta,z) = (r\cos\theta, \, r\sin\theta, \, z).
\]

\begin{proposition}[Jacobian of Cylindrical Coordinates]
$\det J_T(r,\theta,z) = r$.
\end{proposition}

\begin{proof}
The Jacobian matrix, with columns indexed by $r,\theta,z$ and rows by $x,y,z$, is
\[
J_T(r,\theta,z) = \begin{pmatrix} \cos\theta & -r\sin\theta & 0 \\ \sin\theta & r\cos\theta & 0 \\ 0 & 0 & 1 \end{pmatrix}.
\]

Expanding the determinant along the third row, which has only the single nonzero entry $1$ in the $(3,3)$ position, gives
\begin{align*}
\det J_T &= 1 \cdot \det \begin{pmatrix} \cos\theta & -r\sin\theta \\ \sin\theta & r\cos\theta \end{pmatrix} = \cos\theta(r\cos\theta) - (-r\sin\theta)(\sin\theta) \\
&= r\cos^2\theta + r\sin^2\theta = r,
\end{align*}
as claimed.
\end{proof}

Hence, since $r \geq 0$ throughout the natural domain of cylindrical coordinates, $|\det J_T| = r$, and the change of variables formula reads
\[
\iiint_{E'} f(x,y,z)\, dV = \iiint_E f(r\cos\theta, r\sin\theta, z) \; r\, dr\, d\theta\, dz.
\]

\begin{example}[Volume of a Solid Cylinder by Cylindrical Coordinates]
Let $E'$ be the solid cylinder $x^2+y^2 \leq R^2$,
\[
\begin{aligned}
0 &\leq z \\
&\leq h.
\end{aligned}
\]
In cylindrical coordinates, $E$ is the box
\[
\begin{aligned}
0 &\leq r \\
&\leq R,
\end{aligned}
\]

\[
\begin{aligned}
0 &\leq \theta \\
&\leq 2\pi,
\end{aligned}
\]

\[
\begin{aligned}
0 &\leq z \\
&\leq h,
\end{aligned}
\]
so
\[
\operatorname{Vol}(E') = \int_0^{2\pi}\int_0^R \int_0^h r\, dz\, dr\, d\theta = 2\pi \cdot h \cdot \int_0^R r\, dr = 2\pi h \cdot \frac{R^2}{2} = \pi R^2 h,
\]
recovering the elementary formula for the volume of a cylinder.
\end{example}

\subsection{Spherical Coordinates}

Spherical coordinates $(\rho,\varphi,\theta)$ describe a point by its distance $\rho \geq 0$ from the origin, the polar angle $\varphi \in [0,\pi]$ measured from the positive $z$-axis, and the azimuthal angle $\theta \in [0,2\pi)$ measured in the $xy$-plane from the positive $x$-axis, via
\[
T(\rho,\varphi,\theta) = \big( \rho\sin\varphi\cos\theta, \; \rho\sin\varphi\sin\theta, \; \rho\cos\varphi \big).
\]

\begin{proposition}[Jacobian of Spherical Coordinates]

\[
\det J_T(\rho,\varphi,\theta) = \rho^2 \sin\varphi.
\]

\end{proposition}

\begin{proof}
Differentiating each component of $T$ with respect to each of $\rho,\varphi,\theta$ gives the Jacobian matrix
\[
\setlength{\arraycolsep}{3pt}
J_T = \begin{pmatrix} \sin\varphi\cos\theta & \rho\cos\varphi\cos\theta & -\rho\sin\varphi\sin\theta \\ \sin\varphi\sin\theta & \rho\cos\varphi\sin\theta & \rho\sin\varphi\cos\theta \\ \cos\varphi & -\rho\sin\varphi & 0 \end{pmatrix}.
\]

Expand the determinant along the third row. The $(3,3)$ entry is $0$, so only the first two entries of that row contribute:
\begin{align*}
\setlength{\arraycolsep}{3pt}
\det J_T &= \cos\varphi \cdot \det\begin{pmatrix} \rho\cos\varphi\cos\theta & -\rho\sin\varphi\sin\theta \\ \rho\cos\varphi\sin\theta & \rho\sin\varphi\cos\theta \end{pmatrix} \\
&\quad -\; (-\rho\sin\varphi)\cdot \det\begin{pmatrix} \sin\varphi\cos\theta & -\rho\sin\varphi\sin\theta \\ \sin\varphi\sin\theta & \rho\sin\varphi\cos\theta \end{pmatrix}.
\end{align*}
The first minor equals
\begin{align*}
&(\rho\cos\varphi\cos\theta)(\rho\sin\varphi\cos\theta) - (-\rho\sin\varphi\sin\theta)(\rho\cos\varphi\sin\theta) \\
&\quad= \rho^2\cos\varphi\sin\varphi\cos^2\theta + \rho^2\sin\varphi\cos\varphi\sin^2\theta = \rho^2\sin\varphi\cos\varphi,
\end{align*}
using $\cos^2\theta+\sin^2\theta=1$; hence the first term is
\[
\cos\varphi \cdot \rho^2\sin\varphi\cos\varphi = \rho^2\sin\varphi\cos^2\varphi.
\]
The second minor equals
\begin{align*}
&(\sin\varphi\cos\theta)(\rho\sin\varphi\cos\theta) - (-\rho\sin\varphi\sin\theta)(\sin\varphi\sin\theta) \\
&\quad= \rho\sin^2\varphi\cos^2\theta + \rho\sin^2\varphi\sin^2\theta = \rho\sin^2\varphi,
\end{align*}
so the second term is
\[
\rho\sin\varphi \cdot \rho\sin^2\varphi = \rho^2\sin^3\varphi.
\]
Adding the two contributions,
\[
\begin{aligned}
\det J_T &= \rho^2\sin\varphi\cos^2\varphi + \rho^2\sin^3\varphi \\
&= \rho^2\sin\varphi\big(\cos^2\varphi+\sin^2\varphi\big) \\
&= \rho^2\sin\varphi,
\end{aligned}
\]
as claimed.
\end{proof}

Since $\varphi \in [0,\pi]$, $\sin\varphi \geq 0$, so $|\det J_T| = \rho^2\sin\varphi$, and the change of variables formula reads
\begin{align*}
\iiint_{E'} f(x,y,z)\, dV &= \iiint_E f\big(\rho\sin\varphi\cos\theta, \rho\sin\varphi\sin\theta, \rho\cos\varphi\big) \\
&\quad\times\; \rho^2\sin\varphi\; d\rho\, d\varphi\, d\theta.
\end{align*}

\begin{example}[Volume of a Ball by Spherical Coordinates]
Let $E'$ be the solid ball $x^2+y^2+z^2 \leq R^2$. In spherical coordinates, $E$ is the box
\[
\begin{aligned}
0 &\leq \rho \\
&\leq R,
\end{aligned}
\]

\[
\begin{aligned}
0 &\leq \varphi \\
&\leq \pi,
\end{aligned}
\]

\[
\begin{aligned}
0 &\leq \theta \\
&\leq 2\pi,
\end{aligned}
\]
so
\begin{align*}
\operatorname{Vol}(E') &= \int_0^{2\pi}\!\!\int_0^{\pi}\!\!\int_0^{R} \rho^2\sin\varphi \; d\rho\, d\varphi\, d\theta \\
&= \left(\int_0^{2\pi} d\theta\right)\left(\int_0^\pi \sin\varphi\, d\varphi\right)\left(\int_0^R \rho^2\, d\rho\right) = 2\pi \cdot 2 \cdot \frac{R^3}{3} = \frac{4}{3}\pi R^3,
\end{align*}
where the three factors separate because the integrand is a product of a function of $\rho$ alone and a function of $\varphi$ alone, and Fubini's Theorem applies on the box $[0,R]\times[0,\pi]\times[0,2\pi]$; this recovers the classical formula for the volume of a ball.
\end{example}

Figure~\ref{fig:ch11-spherical} identifies the radial distance and the two angles used in spherical coordinates.

\begin{figure}[!htbp]
\centering
\begin{tikzpicture}[scale=1.5,>=Stealth]
  \coordinate (O) at (0,0);
  \coordinate (Xax) at (-2.05,-1.35);
  \coordinate (Yax) at (2.85,0);
  \coordinate (Zax) at (0,3.15);
  \coordinate (P) at (1.8,2.05);
  \coordinate (Q) at (1.8,0.48);
  \draw[->] (O) -- (Xax) node[left,font=\scriptsize]{$x$};
  \draw[->] (O) -- (Yax) node[right,font=\scriptsize]{$y$};
  \draw[->] (O) -- (Zax) node[above,font=\scriptsize]{$z$};
  \draw[dashed,gray!75] (O) -- (Q);
  \draw[dashed,gray!75] (Q) -- (P);
  \draw[vecB,thick] (O) -- (P);
  \node[figred,fill=white,inner sep=1pt,font=\scriptsize]
    at (1.18,1.48) {$\rho$};
  \fill (P) circle (1.3pt)
    node[above right,fill=white,inner sep=1pt,font=\scriptsize]{$P=(\rho,\varphi,\theta)$};
  \draw[vecC] (90:0.92) arc[start angle=90,end angle=49,radius=0.92];
  \node[figgreen,fill=white,inner sep=1pt,font=\scriptsize] at (0.33,0.95) {$\varphi$};
  \draw[gray!80,-{Stealth[length=2mm]}] (0:0.72)
    arc[start angle=0,end angle=15,radius=0.72];
  \node[gray!80,below,fill=white,inner sep=1pt,font=\scriptsize]
    at (0.72,0.18) {$\theta$};
  \fill (Q) circle (1pt) node[below right,font=\scriptsize]{$Q$};
\end{tikzpicture}
\caption[Spherical coordinates]{The spherical coordinates $(\rho,\varphi,\theta)$ of a point $P$: $\rho$ is the distance from the origin, $\varphi$ the angle from the positive $z$-axis, and $\theta$ the azimuthal angle, measured in the $xy$-plane, of the projection $Q$ of $P$.}
\label{fig:ch11-spherical}
\end{figure}

\section{Integration in Arbitrary Dimension}

\subsection{Multiple Integrals and Change of Variables}

The theory developed above for $n=2$ and $n=3$ extends, without any new idea, to functions of $n$ real variables; only the bookkeeping becomes heavier.

\begin{definition}[Integral over a Box in $\mathbb{R}^n$]
Consider the box
\[
B = [a_1,b_1]\times\cdots\times[a_n,b_n] \subset \mathbb{R}^n,
\]

Let $f : B \to \mathbb{R}$ be bounded. A partition of $B$, obtained by partitioning each edge, divides $B$ into subboxes $B_{\mathbf{i}}$, indexed by a multi-index $\mathbf{i} = (i_1,\dots,i_n)$, of $n$-dimensional volume $\Delta V_{\mathbf{i}}$. With
\[
M_{\mathbf{i}} = \sup_{B_{\mathbf{i}}} f,
\]

\[
m_{\mathbf{i}} = \inf_{B_{\mathbf{i}}} f,
\]
define
\[
\begin{gathered}
U(f,P) = \sum_{\mathbf{i}} M_{\mathbf{i}}\, \Delta V_{\mathbf{i}} \\
L(f,P) = \sum_{\mathbf{i}} m_{\mathbf{i}}\, \Delta V_{\mathbf{i}};
\end{gathered}
\]
$f$ is integrable on $B$ if
\[
\sup_P L(f,P) = \inf_P U(f,P),
\]
and this common value is denoted
\[
\int_B f\, dV.
\]

\end{definition}

Integrability of continuous functions, and of functions continuous except on a set of $n$-dimensional content zero, follow exactly as before, by the identical uniform continuity arguments applied coordinatewise; the definition of the integral over a general bounded region $D \subset \mathbb{R}^n$ with $\partial D$ of content zero, via zero-extension to an enclosing box, is likewise unchanged.

\begin{theorem}[Fubini's Theorem in $\mathbb{R}^n$]
Let $B_1 \subset \mathbb{R}^{n-1}$ be a box and define
\[
B = B_1 \times [c,d] \subset \mathbb{R}^{n-1}\times\mathbb{R} = \mathbb{R}^n,
\]
Let $f$ be continuous on $B$. Then
\[
\int_B f \, dV = \int_{B_1} \left( \int_c^d f(\mathbf{x}',x_n)\, dx_n \right) dV',
\]
and the inner $(n-1)$-dimensional integral over $B_1$ may itself be evaluated as a further iterated integral, one variable at a time, in any order.
\end{theorem}

\begin{proof}
By induction on $n$. The base case $n=2$ is the rectangular Fubini Theorem established above. For the inductive step, assume the theorem holds in dimension $n-1$; the identical squeeze argument used in the proof of the two-dimensional case, applied with $x_n$ playing the role of the single distinguished ``outer'' variable and $B_1$ playing the role of the ``inner'' rectangle -- with sums over subrectangles of $B_1$ replacing sums over the index $i$, and Darboux sums over $B_1$-partitions replacing the one-dimensional Darboux sums of that proof -- shows that
\[
A(x_n) := \int_{B_1} f(\mathbf{x}',x_n)\, dV'
\]
is well defined and continuous, and that
\[
\int_B f\, dV = \int_c^d A(x_n)\, dx_n.
\]

By the inductive hypothesis, applied to the continuous function $\mathbf{x}' \mapsto f(\mathbf{x}',x_n)$ on $B_1 \subset \mathbb{R}^{n-1}$ for each fixed $x_n$, the integral $A(x_n)$ may be evaluated as a further iterated integral over the $n-1$ remaining variables, one at a time, in any order; combined with the outer integration over $x_n$, this yields the theorem for dimension $n$, completing the induction.
\end{proof}

\begin{theorem}[Change of Variables Formula in $\mathbb{R}^n$]
Let $T : U \to V$ be a bijective $C^1$ map between open sets $U,V \subset \mathbb{R}^n$ with $\det J_T \neq 0$ on $U$, let $D \subset U$ with $D' = T(D) \subset V$, and let $f$ be continuous on $D'$. Then
\[
\int_{D'} f\, dV = \int_D f\big(T(\mathbf{u})\big) \, |\det J_T(\mathbf{u})| \; dV_{\mathbf{u}}.
\]

\end{theorem}

\begin{proof}[Derivation]
Exactly as in the two- and three-dimensional cases, partition $D$ into small $n$-dimensional boxes and approximate the image, under $T$, of each box by the parallelotope spanned by the $n$ columns of $J_T$, each scaled by the corresponding side length; the $n$-dimensional volume of this parallelotope is the absolute value of the determinant of the matrix formed by its $n$ spanning vectors, that is, $|\det J_T|$ times the volume of the original box, by the same multilinearity argument used above. Passing to the limit of the associated Riemann sums, exactly as before, yields the formula; we omit the details, which repeat those of the two-dimensional derivation verbatim with $n$ coordinates in place of two.
\end{proof}

\subsection{Radial Integrals and Volumes of Balls}

\begin{lemma}[Spherical Shell Formula in $\mathbb{R}^n$]
Let $g : [0,\infty) \to \mathbb{R}$ be continuous, and let $S_{n-1}$ denote the $(n-1)$-dimensional measure of the unit sphere
\[
\{\mathbf{x} \in \mathbb{R}^n : \|\mathbf{x}\|=1\}.
\]
Then
\[
\int_{\mathbb{R}^n} g(\|\mathbf{x}\|)\, dV = S_{n-1} \int_0^{\infty} g(\rho)\, \rho^{n-1}\, d\rho.
\]

\end{lemma}

\begin{proof}[Justification]
This is the direct $n$-dimensional generalization of polar coordinates ($n=2$, where the Jacobian factor $\rho^{n-1}=\rho$ matches the factor $r$ derived above) and spherical coordinates ($n=3$, where $\rho^{n-1}=\rho^2$ matches the factor $\rho^2\sin\varphi$ derived above upon integrating out the angular variable $\varphi,\theta$ over the full unit sphere, which by definition contributes precisely $S_2 = 4\pi$). In general, let $\boldsymbol{\omega}$ range over the unit sphere $S^{n-1}$ and use the radial representation
\[
\begin{aligned}
\mathbf{x} &= \rho\boldsymbol{\omega}, \\
\rho &= \|\mathbf{x}\| \geq 0.
\end{aligned}
\]

The change of variables formula applied to this generalized polar map yields a Jacobian factor depending on $\rho$ alone as $\rho^{n-1}$, exactly as in the cases $n=2,3$ verified explicitly above, times a factor depending only on the angular coordinates of $\boldsymbol\omega$; integrating $g(\rho)\rho^{n-1}$ against this angular factor over all angular coordinates, for each fixed $\rho$, yields by definition the total measure $S_{n-1}$ of the unit sphere, giving the stated formula.
\end{proof}

\begin{example}[Volume of the Unit Ball in $\mathbb{R}^n$]
The volume
\[
V_n = \operatorname{Vol}(\{\mathbf{x}\in\mathbb{R}^n : \|\mathbf{x}\|\leq 1\})
\]
of the unit ball in $\mathbb{R}^n$ can be computed in closed form using the Gaussian integral together with the shell formula above. By Fubini's Theorem in $\mathbb{R}^n$, the $n$-fold Gaussian integral separates into a product of $n$ one-dimensional Gaussian integrals, each equal to $\sqrt{\pi}$ by the earlier example:
\[
\int_{\mathbb{R}^n} e^{-\|\mathbf{x}\|^2}\, dV = \left( \int_{-\infty}^{\infty} e^{-t^2}\, dt \right)^n = \pi^{n/2}.
\]
On the other hand, applying the shell formula with $g(\rho) = e^{-\rho^2}$,
\[
\int_{\mathbb{R}^n} e^{-\|\mathbf{x}\|^2}\, dV = S_{n-1} \int_0^{\infty} e^{-\rho^2}\, \rho^{n-1}\, d\rho.
\]
Substituting $t = \rho^2$, $dt = 2\rho\, d\rho$, so that
\[
\rho^{n-1}\, d\rho = \rho^{n-2}\cdot \rho\, d\rho = t^{(n-2)/2} \cdot (1/2)\, dt,
\]

\[
\int_0^{\infty} e^{-\rho^2}\rho^{n-1}\, d\rho = \frac{1}{2}\int_0^{\infty} e^{-t}\, t^{n/2-1}\, dt = \frac{1}{2}\Gamma\!\left(\frac{n}{2}\right),
\]
by definition of the Gamma function
\[
\Gamma(s) = \int_0^\infty e^{-t}t^{s-1}\,dt.
\]
Equating the two expressions for
\[
\int_{\mathbb{R}^n} e^{-\|\mathbf{x}\|^2}\, dV
\]
and solving,
\[
\pi^{n/2} = S_{n-1}\cdot \frac{1}{2}\Gamma\!\left(\frac{n}{2}\right) \qquad\Longrightarrow\qquad S_{n-1} = \frac{2\pi^{n/2}}{\Gamma(n/2)}.
\]
Finally, applying the shell formula with $g \equiv 1$ but restricted to $\rho \in [0,1]$,
\[
\begin{aligned}
V_n &= \int_0^1 S_{n-1}\, \rho^{n-1}\, d\rho
= S_{n-1} \left[\frac{\rho^n}{n}\right]_0^1 \\
&= \frac{S_{n-1}}{n}
= \frac{2\pi^{n/2}}{n\,\Gamma(n/2)}
= \frac{\pi^{n/2}}{\Gamma\!\left(\frac{n}{2}+1\right)},
\end{aligned}
\]
using the recursive property
\[
\Gamma(s+1) = s\,\Gamma(s)
\]
of the Gamma function with $s=n/2$. As a consistency check, for $n=2$ this gives
\[
\begin{aligned}
V_2 &= \pi/\Gamma(2) \\
&= \pi/1 \\
&= \pi,
\end{aligned}
\]
the area of the unit disk; for $n=3$, using
\[
\begin{aligned}
\Gamma(5/2) &= (3/2)\cdot(1/2)\cdot\sqrt{\pi} \\
&= (3/4)\sqrt{\pi},
\end{aligned}
\]
one obtains
\[
\begin{aligned}
V_3 &= \pi^{3/2}/((3/4)\sqrt{\pi}) \\
&= (4/3)\pi,
\end{aligned}
\]
in agreement with the direct spherical-coordinates computation of the preceding section.
\end{example}

\section{Integral Theorems of Vector Calculus}
\label{sec:vector-integral-theorems}

The gradient, divergence, and curl introduced in Section~\ref{sec:vector-operators} describe local changes in scalar and vector fields. Recall that $\nabla\phi$ is a vector field, $\nabla\cdot\mathbf{F}$ is a scalar field, and $\nabla\times\mathbf{F}$ is a vector field. Positive divergence indicates a local source, while negative divergence indicates a local sink. The curl describes local circulation: its direction is the axis selected by the right-hand rule, and, for a velocity field, its magnitude is twice the local angular speed of the rotational part of the flow. The following theorems make these interpretations precise by connecting local derivatives to global integrals over boundaries.

\subsection{The Divergence Theorem}

\begin{definition}[Flux through a Closed Surface]
Let $\Omega\subset\R^3$ be a bounded domain with piecewise smooth boundary $\partial\Omega$, and let $\mathbf{n}(\bx)$ denote the outward unit normal. The \emph{flux} of a continuous vector field $\mathbf{F}$ through $\partial\Omega$ is
\begin{equation}
\Phi := \iint_{\partial\Omega}\mathbf{F}\cdot\mathbf{n}\,\dd S,
\end{equation}
where $\dd S$ is the surface area element. On a regular parametrized surface $\mathbf{r}(s,t)$, this element is
\[
\dd S=|\mathbf{r}_s\times\mathbf{r}_t|\,\dd s\,\dd t;
\]
the sign of the unit normal is chosen to agree with the prescribed orientation.
\end{definition}

\begin{theorem}[Divergence Theorem, Gauss--Ostrogradsky]
\label{thm:divergence}
Let $\Omega\subset\R^3$ be a bounded domain with piecewise smooth boundary, and let $\mathbf{F}$ be continuously differentiable on a neighborhood of $\overline\Omega$. Then
\begin{equation}
\iiint_\Omega\nabla\cdot\mathbf{F}\,\dd V
=\iint_{\partial\Omega}\mathbf{F}\cdot\mathbf{n}\,\dd S.
\end{equation}

\end{theorem}

\begin{proof}[Proof for Simple Domains]
By linearity, it is enough to establish the identity for each component. For the third component, we must show that
\begin{equation}
\iiint_\Omega\frac{\partial F_3}{\partial z}\,\dd V
=\iint_{\partial\Omega}F_3n_3\,\dd S.
\end{equation}
Assume that $\Omega$ is $z$-simple, with the description
\[
\Omega=\{(x,y,z):(x,y)\in D,\ g_1(x,y)\leq z\leq g_2(x,y)\}.
\]
Integrating first in $z$ and applying the fundamental theorem of calculus gives
\begin{align}
\iiint_\Omega\frac{\partial F_3}{\partial z}\,\dd V
&=\iint_D\left[\int_{g_1(x,y)}^{g_2(x,y)}\frac{\partial F_3}{\partial z}\,\dd z\right]\dd x\,\dd y\notag\\
&=\iint_D\big[F_3(x,y,g_2(x,y))-F_3(x,y,g_1(x,y))\big]\,\dd x\,\dd y.
\end{align}

The boundary consists of the upper graph, the lower graph, and possibly vertical sides. On the upper graph, $n_3\,\dd S=\dd x\,\dd y$; on the lower graph, $n_3\,\dd S=-\dd x\,\dd y$. These signs reflect the outward orientation. On the vertical sides, the normal is horizontal, so $n_3=0$. Therefore
\begin{align}
\iint_{\partial\Omega}F_3n_3\,\dd S
&=\iint_D F_3(x,y,g_2(x,y))\,\dd x\,\dd y\notag\\
&\quad-\iint_D F_3(x,y,g_1(x,y))\,\dd x\,\dd y,
\end{align}
which is exactly the volume integral above. The other components follow by exchanging the roles of $x,y,z$ on domains simple in the corresponding directions. Summing proves the theorem for domains simple in all three directions. For domains decomposable into finitely many such pieces, the fluxes across internal faces cancel because their outward normals are opposite. The extension to general piecewise smooth domains uses localization and approximation, which we omit.
\end{proof}

\begin{remark}[A Multidimensional Fundamental Theorem of Calculus]
The one-dimensional identity
\[
\int_a^b g'(x)\,\dd x=g(b)-g(a)
\]
expresses an integral of a derivative in terms of boundary values. The signs at $a$ and $b$ are precisely the outward normals $-1$ and $+1$. The divergence theorem has the same structure in three dimensions: the volume integral of a derivative becomes a boundary integral weighted by the outward normal. This explains why its proof closely follows the fundamental theorem of calculus.
\end{remark}

\begin{example}[The Continuity Equation]
\label{ex:continuity-equation}
Let $\rho(\bx,t)$ be the density of a conserved quantity, such as mass, charge, or probability, and let $\mathbf{J}(\bx,t)$ be its associated current. Conservation in a fixed volume $\Omega$ means that the rate of change of the total amount equals minus the net outward flux:
\begin{equation}
\frac{\dd}{\dd t}\iiint_\Omega\rho\,\dd V
=-\iint_{\partial\Omega}\mathbf{J}\cdot\mathbf{n}\,\dd S.
\end{equation}

Assuming sufficient regularity to differentiate under the integral and apply Theorem~\ref{thm:divergence}, we obtain
\begin{equation}
\iiint_\Omega\left(\frac{\partial\rho}{\partial t}+\nabla\cdot\mathbf{J}\right)\dd V=0.
\end{equation}

This identity holds for every admissible volume. A continuous integrand whose integral vanishes on every such volume must vanish identically: otherwise, continuity would give a neighborhood on which it has a strict, constant sign, contradicting the integral identity. Hence
\begin{equation}
\frac{\partial\rho}{\partial t}+\nabla\cdot\mathbf{J}=0,
\end{equation}
the \emph{continuity equation}.
\end{example}

Figure~\ref{fig:atlas-outward-flux} connects positive divergence inside a region with outward boundary flux.

\begin{figure}[!htbp]
\centering
\begin{tikzpicture}[scale=1.35,>=Stealth]
\fill[figblue!8] (0,0) circle (1.5);
\draw[figblue,thick] (0,0) circle (1.5);
\foreach \angle in {0,30,...,330} {
  \draw[->,figred,thick] (\angle:1.5) -- (\angle:2.05);
  \draw[->,black!40] (\angle:0.55) -- (\angle:0.95);
}
\node at (0,0) {$\Omega$};
\node[above right,figblue] at (45:1.5) {$\partial\Omega$};
\node[right,figred] at (0:2.05) {$n$};
\end{tikzpicture}
\caption{A planar slice through the unit ball for the three-dimensional field $F(x,y,z)=(x,y,z)$. The field has divergence $3$ and agrees with the outward unit normal on the unit sphere. The resulting flux $4\pi$ equals divergence times volume. Arrows show radial directions, not a surface parametrization.}
\label{fig:atlas-outward-flux}
\end{figure}

\subsection{Green's Theorem in the Plane}
\label{sec:green-plane}

We establish the planar theorem before Stokes' theorem so that the latter can be derived without a circular argument.

\begin{definition}[Line Integrals and Circulation]
For a piecewise smooth oriented curve $\gamma$ with parametrization $\mathbf{r}:[a,b]\to\R^3$, the line integral of a continuous vector field is
\begin{equation}
\int_\gamma\mathbf{F}\cdot\dd\mathbf{r}
=\int_a^b\mathbf{F}(\mathbf{r}(t))\cdot\mathbf{r}'(t)\,\dd t.
\end{equation}
For a closed curve, this is called the \emph{circulation} and is denoted by
\[
\oint_\gamma\mathbf{F}\cdot\dd\mathbf{r}.
\]
In the plane, with $\mathbf{F}=(P,Q)$, the integrand is $P\,\dd x+Q\,\dd y$.
\end{definition}

\begin{theorem}[Green's Theorem]
\label{thm:green-plane}
Let $D\subset\R^2$ be a bounded domain whose boundary is a piecewise smooth simple closed curve, oriented counterclockwise. If $P,Q$ are continuously differentiable on a neighborhood of $\overline D$, then
\begin{equation}
\oint_{\partial D}(P\,\dd x+Q\,\dd y)
=\iint_D\left(\frac{\partial Q}{\partial x}-\frac{\partial P}{\partial y}\right)\dd x\,\dd y.
\end{equation}

\end{theorem}

\begin{proof}[Proof for Simple Domains]
It is enough to prove the identities for $P\,\dd x$ and $Q\,\dd y$ separately. For a $y$-simple domain
\[
D=\{(x,y):a\leq x\leq b,\ h_1(x)\leq y\leq h_2(x)\},
\]
the fundamental theorem of calculus yields
\begin{equation}
\iint_D\frac{\partial P}{\partial y}\,\dd x\,\dd y
=\int_a^b\big[P(x,h_2(x))-P(x,h_1(x))\big]\,\dd x.
\end{equation}

Counterclockwise orientation traverses the lower boundary from $a$ to $b$ and the upper boundary from $b$ to $a$. Any vertical segments contribute nothing because $\dd x=0$. Consequently,
\begin{align}
\oint_{\partial D}P\,\dd x
&=\int_a^b P(x,h_1(x))\,\dd x-\int_a^b P(x,h_2(x))\,\dd x\notag\\
&=-\iint_D\frac{\partial P}{\partial y}\,\dd x\,\dd y.
\end{align}
Similarly, integration first in $x$ on an $x$-simple domain gives
\[
\oint_{\partial D}Q\,\dd y=\iint_D\frac{\partial Q}{\partial x}\,\dd x\,\dd y.
\]

Adding proves the formula when the domain admits both descriptions. Subdivision and cancellation of internal boundary integrals extend the argument to unions of such pieces; the general piecewise smooth case follows by approximation.
\end{proof}

\begin{remark}[The Divergence Form of Green's Theorem]
Applying the circulation formula to the rotated field $(-Q,P)$ gives
\begin{equation}
\oint_{\partial D}(P\,\dd y-Q\,\dd x)
=\iint_D\left(\frac{\partial P}{\partial x}+\frac{\partial Q}{\partial y}\right)\dd x\,\dd y.
\end{equation}
For counterclockwise orientation,
\[
\mathbf{n}\,\dd s=(\dd y,-\dd x),
\]
so the left-hand side is the outward flux of $(P,Q)$. This is the two-dimensional divergence theorem: in the plane, circulation and divergence are related by a rotation through $90^\circ$.
\end{remark}

\subsection{Stokes' Theorem}

Let $\Sigma\subset\R^3$ be an oriented piecewise smooth surface with boundary $\partial\Sigma$. The boundary orientation is induced by the right-hand rule: when the right thumb points along the chosen normal $\mathbf{n}$, the fingers indicate the positive direction around the surface.

\begin{theorem}[Kelvin--Stokes Theorem]
\label{thm:stokes}
If $\mathbf{F}$ is continuously differentiable on a neighborhood of $\Sigma$, then
\begin{equation}
\iint_\Sigma(\nabla\times\mathbf{F})\cdot\mathbf{n}\,\dd S
=\oint_{\partial\Sigma}\mathbf{F}\cdot\dd\mathbf{r}.
\end{equation}

\end{theorem}

\begin{proof}[Proof Sketch]
First suppose that $\Sigma$ is the graph
\[
\mathbf{r}(x,y)=(x,y,g(x,y))
\]
over a planar domain $D$, oriented upward, with $g$ sufficiently smooth. Writing $\mathbf{F}=(F_1,F_2,F_3)$ and evaluating its components on the graph gives
\[
\mathbf{F}\cdot\dd\mathbf{r}
=(F_1+F_3g_x)\,\dd x+(F_2+F_3g_y)\,\dd y.
\]
Set
\[
\begin{gathered}
P=F_1\circ\mathbf{r}+(F_3\circ\mathbf{r})g_x \\
Q=F_2\circ\mathbf{r}+(F_3\circ\mathbf{r})g_y.
\end{gathered}
\]
The chain rule and cancellation of mixed derivatives show that
\[
Q_x-P_y=(\nabla\times\mathbf{F})(\mathbf{r}(x,y))\cdot(-g_x,-g_y,1).
\]
The oriented surface element is
\[
\mathbf{n}\,\dd S=(-g_x,-g_y,1)\,\dd x\,\dd y,
\]

Green's theorem gives the desired equality. A general surface is treated by subdividing it into local graph patches. Adjacent patches traverse their common boundary in opposite directions, so internal contributions cancel, leaving only $\partial\Sigma$.
\end{proof}

\begin{remark}[A Common Boundary Principle]
Stokes' theorem again has the form ``integral of a derivative over a region equals integral of the original field over its boundary.'' Here the derivative is the curl, the region is a two-dimensional surface, and its boundary is a one-dimensional curve. Green's circulation theorem is the special case $\Sigma=D$ in the $xy$-plane, $\mathbf{n}=(0,0,1)$, and $\mathbf{F}=(P,Q,0)$, because $(\nabla\times\mathbf{F})\cdot\mathbf{n}=Q_x-P_y$. Thus the two forms of Green's theorem connect the planar versions of both Stokes' and Gauss' theorems.
\end{remark}
